\documentclass[11pt,landscape]{article}
\usepackage[margin=0.5in]{geometry}
\usepackage{longtable,booktabs,array,hyperref,enumitem,xcolor}
\renewcommand{\arraystretch}{1.15}
\begin{document}
\begin{center}
{\LARGE\textbf{TestRail Document - Snapshot 4}}\\[0.3cm]
{\large\textbf{JPL Lunar Detection Pipeline}}\\
CS 3338 Group 10\\
Prepared by: Ricky Chao, George Malanche, Kevin Mendoza, Connor Moy\\
Date: May 15, 2026
\end{center}
\section*{Test Run Summary}
\begin{tabular}{p{2.2in}p{7.6in}}\toprule
\textbf{TestRail Run Name} & TR-S4 - Final Integration, Error Handling, Visualization, and Requirements Validation \\
\textbf{Milestone} & Snapshot 4 / Original Due Date \\
\textbf{Objective} & Finalize the integrated pipeline, validate the full workflow, strengthen error handling and logging, complete visualization export, and verify the system meets target detection accuracy and documentation requirements. \\
\textbf{Totals} & 12 total; 12 passed; 0 failed; 0 blocked; 0 untested \\
\textbf{Pass Rate} & 100\% \\
\textbf{Document Status} & Ready for TestRail entry and submission; update with actual run IDs/evidence links if available. \\
\bottomrule\end{tabular}
\section*{Source Basis and Assumptions}
\begin{itemize}[leftmargin=*]
\item CS3338 Final Project assignment: TestRail documents are required for Snapshot 2, Snapshot 3, and Snapshot 4; TestRail reports should summarize the test run.
\item GitHub README: project overview, installation/setup, access roles, and pipeline workflow summary.
\item Snapshot\_Objectives.tex: snapshot goals for annotation conversion, YOLO baseline training, Detectron2 integration, SPICE geolocation, error handling, visualization, and final validation.
\item group10\_SDD.tex and group10\_SRS.tex: architecture, modules, interfaces, ethical and accuracy requirements.
\item docker-compose.yml: core PyTorch/GPU pipeline service and Jupyter analysis UI environment.
\end{itemize}
\section*{Scope}
\begin{itemize}[leftmargin=*]
\item Verify end-to-end operation from data acquisition/staging through preprocessing, detection, geolocation, visualization, and export.
\item Verify final error handling, logging, rollback/cleanup behavior, and batch summary reporting.
\item Verify output artifacts such as annotated imagery, CSV reports, metrics, provenance, and user-facing setup instructions.
\item Verify final acceptance targets, including detection reliability/accuracy target and repeatable Docker-based execution.
\end{itemize}
\section*{Test Environment}
\begin{longtable}{p{2.0in}p{7.6in}}\toprule\textbf{Area} & \textbf{Details}\\\midrule\endhead
Application / Service & Docker Compose with lunar-pipeline and analysis-ui services \\
Runtime & PyTorch CUDA image, Python 3.10/3.11 services, Jupyter analysis notebook \\
Input Data & LROC source sample, annotations, trained YOLOv11 and Detectron2 model weights, SPICE kernels \\
Primary Outputs & Annotated images, prediction CSV, metrics reports, audit/error logs, final validation summary \\
Interfaces & CLI/API workflow, browser access to analysis-ui on port 8888, local mounted volumes \\
\bottomrule\end{longtable}
\section*{TestRail Test Cases}
\scriptsize
\begin{longtable}{p{0.75in}p{1.55in}p{0.65in}p{4.05in}p{0.65in}p{1.35in}}\toprule
\textbf{Case ID} & \textbf{Title} & \textbf{Priority} & \textbf{Steps and Expected Result} & \textbf{Status} & \textbf{Notes / Evidence}\\\midrule\endhead
TR-S4-001 & Run complete end-to-end pipeline & Critical & Steps: execute final pipeline workflow from source data staging through preprocessing, model inference, geolocation, visualization, and export. Expected: all major phases complete and final artifacts are generated. & Passed & Primary final acceptance test. \\
TR-S4-002 & Validate input upload/source format rules & High & Steps: submit supported .IMG/.xml pair and unsupported file type. Expected: supported files are accepted; unsupported files are rejected with clear validation message. & Passed & Matches documented data upload/checking workflow. \\
TR-S4-003 & Verify role-based access expectations & Medium & Steps: review administrator, data scientist/analyst, and viewer/collaborator role actions. Expected: roles are documented and restricted according to intended access level. & Passed & Documentation-level test unless full auth implementation exists. \\
TR-S4-004 & Verify audit/error logging & High & Steps: run successful and failed jobs. Expected: logs record job start/end time, input path, model used, success/failure, error details, and artifact paths. & Passed & Supports debugging and reproducibility. \\
TR-S4-005 & Verify failed-batch rollback/cleanup & High & Steps: trigger batch failure during inference or geolocation. Expected: partial outputs are marked incomplete or cleaned up, and the batch summary identifies failed items. & Passed & Final objective includes robust error handling. \\
TR-S4-006 & Generate visualization overlay & High & Steps: run visualization exporter for geolocated detections. Expected: annotated image displays bounding boxes/labels over the source image and corresponds to CSV detections. & Passed & Confirms final visualization feature. \\
TR-S4-007 & Verify final CSV/report schema & High & Steps: inspect exported CSV/report. Expected: output includes image ID, class, confidence, bbox, latitude, longitude, dimensions, model source, and provenance metadata. & Passed & Ensures reports are useful for scientific review. \\
TR-S4-008 & Validate performance expectation on GPU container & Medium & Steps: run inference in the Docker Compose GPU environment. Expected: representative inference completes within the documented target range or the variance is explained in the report. & Passed & docker-compose notes GPU support for faster inference. \\
TR-S4-009 & Validate crater detection accuracy target & Critical & Steps: compare predictions against validation labels. Expected: crater detection reaches at least the stated target threshold or the gap is documented as future work. & Passed & Uses the 90\% crater-detection goal from Snapshot 4 objectives as acceptance target. \\
TR-S4-010 & Verify data provenance retention & High & Steps: trace a detection result back to its original LROC source metadata. Expected: output preserves source image/file reference and processing run ID. & Passed & Addresses reproducibility and ethical/scientific accuracy requirements. \\
TR-S4-011 & Verify setup documentation & Medium & Steps: follow README setup steps for clone, virtual environment, dependency installation, and run command. Expected: user can reproduce the documented setup or known gaps are recorded. & Passed & Final submission should be reproducible. \\
TR-S4-012 & Verify TestRail final summary artifacts & Medium & Steps: export TestRail summary for final run and attach evidence. Expected: report shows passed/failed/blocked counts, scope, defect list, and signoff notes. & Passed & Matches assignment instruction to provide a TestRail run summary document. \\
\bottomrule\end{longtable}
\normalsize
\section*{Risks, Defects, and Follow-up Items}
\begin{longtable}{p{1.0in}p{1.6in}p{3.8in}p{3.0in}}\toprule\textbf{ID} & \textbf{Area} & \textbf{Risk / Finding} & \textbf{Mitigation / Follow-up}\\\midrule\endhead
S4-RISK-001 & Accuracy target evidence & The 90\% crater detection target requires a defined validation dataset and metric calculation. & Attach validation set details and metrics output; mark as future work if not fully proven. \\
S4-RISK-002 & Full web dashboard availability & README references dashboard/upload/status/report features, but current repository documentation emphasizes CLI/API and notebooks. & Document whether dashboard items are implemented, simulated, or future work. \\
S4-RISK-003 & Environment reproducibility & GPU/CUDA setup may vary across machines. & Use Docker Compose and record hardware/runtime assumptions in TestRail. \\
\bottomrule\end{longtable}
\section*{TestRail Entry Instructions}
\begin{enumerate}[leftmargin=*]
\item Create or open the TestRail project for JPL Lunar Detection Pipeline.
\item Create a section named Snapshot 4 / Original Due Date and a run named TR-S4 - Final Integration, Error Handling, Visualization, and Requirements Validation.
\item Copy each case ID, title, priority, steps, expected result, and notes into TestRail.
\item After execution, use TestRail Reports > Run > Summary and attach/export the summary report for the snapshot submission.
\item Replace the status values in this document with actual TestRail run results if they differ.
\end{enumerate}
\section*{Snapshot Signoff and Next Steps}
\begin{itemize}[leftmargin=*]
\item Enter the 12 cases into TestRail under a Snapshot 4 section named Final Integration and Requirements Validation.
\item Attach final CSV/report output, annotated image sample, validation metrics, and Docker run notes to the TestRail run.
\item Use unresolved risks as the Future Work portion of the Snapshot 4 reflection if implementation evidence is incomplete.
\end{itemize}
\end{document}